\label{sec:conclusion}

The F.3 resolution rule, applied to projected 2030 seat counts, identifies 43 states
requiring block-group adjacency for the 2030 redistricting cycle. The block-group
adjacency graphs are 5$\times$ larger than tract-level graphs, require 21 CPU-hours
per state to construct, and take 2 weeks to validate. If constructed during the
redistricting window, they consume 4--6 weeks of the available 6--12 months.
If pre-built by January 2031, they reduce the redistricting preparation overhead
from months to days.

The recommendation is unambiguous: begin block-group adjacency construction for
the 43 states no later than January 2030. The construction can be parallelized
across states on any modern multi-core server, and the total resource cost (3--4 weeks
of elapsed time, 4 GB of storage, and 64 GB peak RAM) is modest compared to the
operational value of having the infrastructure ready when census data arrives.

\subsection{Implementation Steps}

The recommended implementation sequence for the Q.3 pre-build:

\begin{enumerate}
\item \textbf{January 2029}: Identify the 43 states requiring block-group resolution,
  using Q.4 projected seat counts. Confirm the list against Census Bureau 2030
  block-group geography when available (expected release: 2029).

\item \textbf{July 2029}: Test the \texttt{bisect fetch --resolution block-group}
  extension on a single representative state (Texas) using 2020 block-group geometry
  as a proxy for 2030. Validate construction process and resource estimates.

\item \textbf{January 2030}: Begin full 43-state block-group adjacency construction.
  Run 4 states simultaneously on a 64 GB machine. Target completion: April 2030
  (13 weeks elapsed).

\item \textbf{April--June 2030}: Run validation protocol for all 43 states.
  Address any connectivity failures or boundary anomalies. Final validation report
  by June 2030.

\item \textbf{July 2030}: Deliver validated block-group adjacency graphs to the
  pre-registration configuration as part of the Q.2 algorithm pre-registration.

\item \textbf{January 2031}: Final infrastructure check. Verify that 2030 block-group
  shapefiles (when available from the Census Bureau) match the 2029-vintage geometry
  used for pre-construction. If changes exist (rare between decennial and pre-release
  TIGER), run differential updates for affected states.
\end{enumerate}

\subsection{Risk Mitigation}

The primary risk in the Q.3 pre-build is that the 2030 block-group geography differs
substantially from the 2020 geometry used for pre-construction. The Census Bureau
updates TIGER/Line geometry between census years, and block-group boundaries can change.
Mitigation: run the pre-build using 2030 TIGER/Line shapefiles as soon as they become
available (expected: 2029), not the 2020 shapefiles. If 2030 shapefiles are not
available by January 2030, run using 2020 shapefiles and schedule a differential
update once 2030 shapefiles are released.

\bigskip
\noindent\textit{The author thanks the TIGER/Line program staff at the Census Bureau
for public access to block-group geometry shapefiles and the \textsc{bisect-data}
development team for the adjacency construction infrastructure. All errors are the author's own.}
